\FloatBarrier
\section{Experimental Protocol}
\label{app:protocol}
\subsection{Inference configuration}
The principal checkpoint is the released \texttt{Qwen/Qwen3-8B}, revision \texttt{b968826d9c46}. It has 36 layers, width 4096, 32 query heads, eight KV heads, and head width 128. Inference uses bf16 and SDPA. The native maximum position is 40,960, with RoPE base $10^6$ and no active scaling unless a YaRN arm is explicitly named. The study does not continue-pretrain the backbone.

\paragraph{Split-depth configuration.}
\label{app:split-choice}
The released model registry sets \texttt{DEFAULT\_RATIO=0.33} and maps Qwen3-8B explicitly to $j=12$; for an unlisted backbone, \texttt{--j auto} falls back to $\operatorname{round}(0.33L)$. Training and timing entry points also default to 12, while an explicit $j$ overrides the default. This is a deployment heuristic motivated by the depth--fidelity trade-off; the study does not establish it through probe-based or held-out benchmark selection. The main Qwen3-8B evaluations retain $j=12$ across tasks and source lengths, including the no-LoRA control; Table~\ref{tab:multidepth} reports sensitivity to alternative splits. Its separately trained adapters and varying parameter counts prevent interpreting the sweep as an equal-training-compute optimization or establishing a universally optimal depth.

The default writer independently encodes disjoint 512-token chunks through blocks 0--11, restarting positions at zero. It stores all target-token residuals in bf16. The reader orders selected chunks by document position, prepends the BOS sink, and appends the independently encoded query. The nominal cap is $1+12\times512+512=6{,}657$ tokens; shorter tails and queries produce observed means around 6.2--6.5k. Its causal mask allows later chunks to attend to earlier chunks. In the block-diagonal control, each chunk sees itself and the sink, while the final query still sees every chunk.

BM25 operates on token-ID documents and question token IDs, using $k_1=1.5$, $b=0.75$, and Robertson IDF with $+1$ smoothing. There is no lowercasing, stemming, or stop-word removal. The saved \texttt{rounds=0} setting means automatic $\lceil k/h\rceil$ rounds, not a single pass. With $k=12$ and hop width $h=4$, all five benchmark families use three rounds, with newly selected chunks forming the next frontier. Decoding is greedy without a chat template or thinking mode. The MidCache and MidCache without LoRA configurations both use $j=12$; Appendix~\ref{app:frozen-recheck} specifies the no-adapter evaluation protocol.

\paragraph{Persistence and query boundaries.}
\label{app:query-boundary}
The reusable interface accepts source and query separately: source tokenization and chunk boundaries are fixed before questions arrive. A chunk ID addresses both its token-ID retrieval metadata and its bf16 residual matrix. Changing only the upper-layer adapter or the selected IDs preserves the stored lower-layer states. Changing source tokens, chunk boundaries, split depth, lower-layer weights, or Write positions requires rebuilding the affected states; overlap also makes a chunk depend on its left prefix. The storage formula counts residual payload only, excluding token IDs, retrieval-index metadata, and request-local KV.

The released monolithic benchmark driver instead tokenizes the complete formatted prompt and takes its last chunk as $q$, with preceding chunks as candidate evidence. BM25 receives the bare question. Thus $q$ need not coincide with the question text: it can contain a document tail, and a question can straddle a chunk boundary. Any block containing query-dependent text must be rebuilt for that query. These formatted-example quality tests do not themselves measure cross-query reuse. The source/query-separated serving experiments measure reuse with fixed document blocks. All methods within each paired control share the same boundary.

The benchmark sink uses the tokenizer's BOS when available and otherwise the first formatted-input token. The Qwen3 tokenizer has no BOS ID. The overlap/KV confirmation and distillation validation instead explicitly use model BOS 151643, identically across their arms; their Qasper values therefore need not equal the main-table Qasper scores. This convention changes an input token, not the chunk selector or the number of sink positions.

\subsection{Suffix-adapter training}
The rank-32 LoRA has $\alpha=32$, zero dropout, and 58.20M trainable parameters (0.71\% of the model). It updates Q/K/V/O and gate/up/down projections only in layers 12--35. The adapter-disabled backbone is the frozen teacher. At every query token, teacher top-64 indices define the shared support used by Equation~\ref{eq:distill}; temperature is one and loss is averaged over query tokens. There is no auxiliary cross-entropy. Appendix~\ref{app:distillation-validation} measures support mass and full-vocabulary divergence on held-out books after training.

Training uses a 64-book PG-19 training subset in 4,096-token windows, partitioned into seven context chunks and one query chunk of 512 tokens each. Teacher and student share the sink, token IDs, chunk order, and query segment. The teacher disables the adapters and processes $[s;x_1;\ldots;x_7;q]$ from layer zero as one causal sequence. Its document positions see the sink and preceding document positions, but no query tokens; its query positions see all preceding document positions and their own causal prefix.

The student's lower layers instead process the sink, each context chunk, and the query as separate causal segments, without gradients or an added per-chunk sink. If $g(t)$ identifies a segment in the assembled order, lower-layer visibility is $u\le t$ and $g(u)=g(t)$; teacher and upper-layer visibility is simply $u\le t$. Write positions restart at zero for each segment. Upper-layer positions are contiguous across the 4,097-position pack, including the sink. Only the suffix adapter receives gradients. This is a causal teacher with broader document context than the student's lower layers, not a teacher with access to future query tokens.

During generation, the lower-layer KV cache contains the query and generated prefix; the upper-layer cache contains the sink, selected document states, query, and generated prefix. A new token uses its next query-local position below the split and its next pack position above it. The continuous-prefix evaluation oracle uses the same full-pack visibility as the teacher but retains the evaluation adapter, whereas training disables it.

AdamW uses $\beta=(0.9,0.95)$, zero weight decay, peak learning rate $10^{-4}$, 50 warmup steps, cosine decay to zero, clipping at 1.0, and seed 42. The backbone is bf16; adapter master parameters are fp32 and matrix products use bf16 autocast. The executed run uses non-reentrant gradient checkpointing, 4,000 steps, one GPU, one 4,096-token window per step, and no accumulation: 16.384M training tokens. The 7.232M-token corpus is traversed cyclically, without inserted document separators; the initial eight windows are skipped once and can recur after wrapping. Training peak allocated memory is 18.01 GiB. These settings follow the checkpoint configuration and executed training record; total compute across preliminary experiments is not reported.

\subsection{Benchmarks and scoring}
\begin{table}[htbp]
\centering
\small
\setlength{\tabcolsep}{4pt}
\begin{tabular}{@{}lll@{}}
\toprule
Benchmark & Support / aggregation & Generation \\
\midrule
RULER & 15 cells, 100/cell; cell macro & 48; VT 60 \\
LongEval & 100/length; five-length macro & 16; without LoRA 48 \\
LongBench & 1,150 examples; six-dataset F1 macro & 32/64/128 \\
BABILong & 21 cells, 100/cell; cell macro & 20 \\
LoCoMo & 1,986 items; item-weighted Judge & 48; scale study 512 \\
\bottomrule
\end{tabular}
\caption{Evaluation supports and maximum generated tokens. RULER uses official case-insensitive string matching; BABILong uses official answer comparison with task labels. LongEval extracts the first numeric answer, LongBench uses official multi-reference token F1, and LoCoMo combines semantic judging with local abstention scoring. RULER, LongBench, and BABILong use seed 42; LongEval uses seed 1234. Adapted runs use eight shards for RULER/LongEval and four for BABILong/LoCoMo. The three paired LongBench configurations and no-adapter runs use four shards throughout.}
\label{tab:eval}
\end{table}

RULER Cohort A uses \texttt{niah\_single\_2}, \texttt{niah\_multikey\_1}, and \texttt{variable\_tracking}. Cohort B uses an independently sampled \texttt{niah\_single\_3} cohort for the matched YaRN and raw-replay comparisons. Both average 15 cells across 8k, 16k, 32k, 64k, and 128k with 100 examples per cell. The MidCache means are 97.05 and 96.07, respectively. A and B are labels for these experimental sample sets, not official RULER variants or train/test splits. The main table uses the task--length support of set A, with independently sampled synthetic runs for the adapter and no-adapter configurations. Set B is reserved for paired replay and YaRN controls. These macros cover 8k--128k; no direct difference is taken between the two cohorts.

\paragraph{Baseline mechanisms and implementation scope.}
The released driver implements the KV-Direct quality reference as full-depth, full-context recomputation without retrieval or LoRA. It evaluates that recomputation endpoint rather than the native system's residual-checkpoint storage and restoration scheduler. The HCache-style overview evaluates a retrieval-free split checkpoint; it is distinct from the same-$j=12$ adapter diagnostic and does not reproduce native HCache. StreamingLLM's sink/window reference retains initial and recent token IDs before dense generation. The SnapKV/PyramidKV controls first prefill the complete prompt and then compress retained KV. MemoryLLM uses its separately released backbone, as marked in Table~\ref{tab:accuracy}. These implementation scopes prevent attributing native serving-system costs to the quality overview. Appendix~\ref{app:chunkkv} separately describes the same-evidence CacheBlend-style reference and its simplified repair policy.

LongEval's common-support mean covers 8k--128k. LongBench averages the six dataset F1 scores equally; the three-task LLoCO mean is kept separate. BABILong averages seven lengths per task and three tasks for the overview. LoCoMo's full-set Judge folds 1,540 answerable questions and 446 locally scored abstention questions into one item-weighted score.

\paragraph{Software and artifacts.}
The benchmark environments use Python 3.10+, PyTorch 2.10, Transformers 5.5.4, PEFT 0.10+, Datasets 2.14+, bf16, and SDPA unless specified otherwise. The three paired LongBench configurations, KV control, and no-adapter benchmark runs use the B300 environment in Appendix~\ref{app:frozen-recheck}; Appendix~\ref{app:systems} gives the device-specific timing environments. Reproduction requires the evaluation code, adapter configuration, saved predictions or permissible score records, judge decisions, and timing records. External data remain subject to their original providers' terms.
